\section{Introduction}
\label{sec:introduction}

Research has historically been constrained by the number, ability, and organization
of human researchers. Artificial intelligence may relax this constraint. The
relevant change would not occur when AI becomes useful in ordinary production, or
even when it automates a large share of existing jobs. It would occur when AI systems
can perform the research tasks required to produce better AI systems. At that point,
AI capability would become an input into its own law of motion, and innovation would
no longer be tied mechanically to the supply of human researchers.

The paper does not model the earlier transition from an economy without commercially
useful AI to one that uses AI in ordinary production. One can think of that
pre-AI period as an economy with conventional production and human research, in
which scientific knowledge, algorithms, and computing capacity accumulate before AI
services become economically useful. A formal model of that transition would require
an adoption margin and an explanation of how research is financed before it produces
current AI revenues. We instead begin after productive AI services already exist but
while frontier research remains primarily dependent on human researchers. The object
of the paper is the subsequent automation of AI research.

This possibility raises an economic question before it raises a forecasting question.
Suppose frontier improvements combine human researchers and automated-research
services. Firms then choose both the scale and the composition of research. When the
two inputs are gross substitutes, the automated share increases if the human research
wage rises relative to the unit resource cost of automated research. Automation is
therefore a continuous equilibrium transition rather than a date imposed on the
model. Human researchers can become asymptotically inessential, but substitution
above one does not guarantee that outcome: the wage must also rise relative to the
cost of automated research.

Population growth is important during the human-dominated phase. In standard R\&D-based growth
models, new ideas depend on a growing but scarce human research input
\citep{romer1990,jones1995,aghionhowitt1992}. With diminishing returns to the stock of
ideas, a constant research-employment share implies semi-endogenous growth: the
long-run rate of innovation is pinned down by the growth of the research labor force.
Our benchmark preserves this logic while human research dominates. If population grows at rate
$n$, the capability growth rate is $\eta n/(1-\phi)$; without population growth,
positive asymptotic growth disappears when $\phi<1$.

Research automation can alter this result. Effective research labor can expand
through additional automated-research services rather than through population growth.
When those services dominate, capability growth is tied to their equilibrium growth
rate. Innovation can accelerate if higher output finances enough additional
automated research, but diminishing returns to ideas and research effort can preserve
a constant-growth path. The outcome depends on the elasticities in production and
research and on the allocation of output to research.

Recent discussions often move too quickly from this mechanism to one of two
conclusions. One view treats an intelligence explosion as the natural implication of
automating AI research. The other assumes that physical bottlenecks and diminishing
returns will necessarily preserve familiar growth rates. Neither conclusion follows
without specifying the ideas production function. Scenario exercises such as
\citet{kokotajlo2025} make the feedback mechanism concrete, but rapid takeoff depends
on several jointly necessary assumptions about research automation, algorithmic
progress, compute, and organizational capacity. An economic model should expose those
assumptions separately.

Market structure determines both current use and the private return to innovation.
Frontier AI requires large fixed investments and may exhibit economies of scale and
scope \citep{korinekvipra2025,gans2024}. A monopolist restricts the current supply of
AI services but captures operating rents that raise the private value of better
capability. We therefore begin with one integrated developer that both supplies
quality-adjusted AI services and invests in the frontier. Competitive final-good
firms combine capital with a CES composite of production labor and AI services. The
elasticity of substitution inside this composite governs adoption and the demand
elasticity faced by the monopolist.

The same elasticity separates the response of the real wage from the response of the
production-labor share. An expansion of AI services reduces this share when AI and
production labor are gross substitutes, but the wage falls only when substitution is
sufficiently strong. There is therefore an intermediate region in which production
workers receive a smaller share of output even though their wage rises. This
distinction is important for the distributional analysis: a declining labor share is
not, by itself, evidence that the real wage has fallen. In the full model, the
aggregate labor share also depends on how workers are reallocated between production
and research.

The resulting welfare problem has two distinct margins. Static market power reduces
the current use of AI services. At the same time, consumer surplus and knowledge
spillovers imply that the developer may capture only part of the social value of
better AI and therefore invest too little in research. Reducing the markup can
correct the deployment distortion while also eroding the rents that finance
innovation. Price regulation and research policy must consequently be studied
together. The present version isolates this market-structure mechanism.

The toy model delivers five results. First, the automated share of research follows
a closed-form law driven by the research-substitution elasticity and the wage relative
to the resource cost of automated research. Second, population growth sustains
semi-endogenous capability growth while human research dominates. Third, when
automated research dominates, a constant-growth path exists only if diminishing
returns offset the feedback between output and research. Fourth, the production CES
separates a labor-bottlenecked regime, a constant-share benchmark, and a regime in
which the return to capital rises with capability. Fifth, the production structure
yields distinct substitution thresholds for the real wage and the
production-labor share.

The extended model retains the integrated monopolist and adds distributional
heterogeneity and richer capital and compute accumulation. It will compare a
laissez-faire monopoly with regulated access prices, research subsidies or prizes,
and policies that combine both instruments. The objective is not to select a policy
before solving the model. It is to identify when lower AI-service prices must be
paired with explicit support for frontier research.

The rest of the paper is organized as follows. Section~\ref{sec:literature} relates
the framework to endogenous growth, transformative AI, automation and labor income,
measurement of AI services, and industrial organization. Section~\ref{sec:toy}
develops the toy model.
Section~\ref{sec:extended} sets out the extended general-equilibrium model.
Section~\ref{sec:regulation} defines the regulatory scenarios.
Section~\ref{sec:conclusion} concludes.
